\documentclass[11pt]{article}
\usepackage[margin=1in]{geometry}
\usepackage[T1]{fontenc}
\usepackage{lmodern,microtype,hyperref,enumitem,amsmath,amssymb,xcolor,booktabs}
\usepackage[numbers,sort&compress]{natbib}
\usepackage{url}
\hypersetup{colorlinks=true,urlcolor=blue,citecolor=blue,linkcolor=blue}

\title{Q11 --- Deeper projects:\\ Why do models hallucinate, and can they know when they are unsure?}
\author{DLS 2026 --- Graduate AI-Internals Module}
\date{}

\begin{document}
\maketitle

\section*{Framing the question}

A language model trained by maximum-likelihood on next-token prediction is
optimized to produce \emph{fluent and plausible} continuations, not
\emph{truthful} ones. Hallucination---confident generation of content that is
unsupported or false---is therefore not an incidental bug but a predictable
consequence of the training objective, the data distribution, and the decoding
procedure. Recent analyses argue that hallucinations arise partly because
pretraining rewards guessing over abstention, and that standard benchmarks and
reward schemes penalize ``I don't know'' responses, incentivizing confident
fabrication \citep{kalai2025hallucinate,maynez2020faithfulness}. A closely
related question is \emph{calibration}: can a model's own probabilities, hidden
states, or verbalized confidence reliably signal when it is likely to be wrong?
Evidence is mixed---base models can be well calibrated on some tasks
\citep{kadavath2022know}, but RLHF can degrade calibration
\citep{openai2023gpt4}, and factual knowledge appears to be linearly encoded in
intermediate activations in ways that sometimes diverge from the model's own
outputs \citep{azaria2023internal,burns2023discovering}.

The projects below ask you to measure these phenomena carefully, probe their
mechanistic origins, and build interventions. Each project lists concrete
deliverables. Use open-weight models (e.g.\ Llama, Mistral, Pythia) so that
internal states are accessible; proprietary APIs may be used for comparison.

\section{Project A --- Measuring and decomposing hallucination}

\textbf{Goal.} Build a reproducible pipeline that elicits, detects, and
categorizes hallucinations, and that separates \emph{knowledge gaps} (the model
never learned the fact) from \emph{retrieval/decoding failures} (the model
``knows'' but outputs something false).

\textbf{Tasks.}
\begin{enumerate}[leftmargin=*]
  \item Assemble an evaluation set combining a closed-book factual QA
    benchmark, a summarization faithfulness benchmark, and an adversarial
    truthfulness benchmark \citep{lin2022truthfulqa}. Annotate each item with a
    ground-truth answer and supporting evidence.
  \item Implement at least two automatic detectors: (i) an entailment/NLI-based
    faithfulness check against source documents
    \citep{maynez2020faithfulness}, and (ii) a sampling-based consistency score
    such as SelfCheckGPT, which flags claims that vary across stochastic
    resamples \citep{manakul2023selfcheckgpt}.
  \item Distinguish knowledge gaps from recall failures: for each item, probe
    whether the correct answer appears among top-$k$ sampled completions or
    under constrained decoding even when greedy decoding is wrong.
\end{enumerate}

\textbf{Deliverables.}
\begin{itemize}[leftmargin=*]
  \item A documented code repository and a dataset card with annotation
    guidelines and inter-annotator agreement on a 100-item subset.
  \item A report quantifying hallucination rate per task and per model size,
    with a confusion-style breakdown into knowledge-gap vs.\ recall-failure
    categories.
  \item An error taxonomy with at least 20 qualitative examples per category.
\end{itemize}

\section{Project B --- Confidence, calibration, and abstention}

\textbf{Goal.} Determine how well a model ``knows when it is unsure,'' comparing
sequence likelihoods, verbalized confidence, and a trained ``P(True)'' head.

\textbf{Tasks.}
\begin{enumerate}[leftmargin=*]
  \item Compute reliability diagrams and expected calibration error (ECE) for
    (i) token/sequence probabilities, (ii) verbalized numeric confidence
    elicited by prompting, and (iii) a self-evaluation probe asking the model
    whether its own answer is correct \citep{kadavath2022know}.
  \item Compare a base/instruct checkpoint against an RLHF-tuned variant to test
    whether alignment tuning degrades calibration \citep{openai2023gpt4}.
  \item Build an abstention policy: choose a confidence threshold that optimizes
    a scoring rule which rewards correct answers, mildly rewards abstention, and
    heavily penalizes confident errors, as advocated for discouraging guessing
    \citep{kalai2025hallucinate}. Plot accuracy-vs-coverage (selective
    prediction) curves.
\end{enumerate}

\textbf{Deliverables.}
\begin{itemize}[leftmargin=*]
  \item Reliability diagrams and ECE/AUROC tables for all confidence signals
    across $\geq 2$ models.
  \item A selective-prediction (risk--coverage) analysis with the derived
    abstention thresholds.
  \item A short discussion of which confidence signal is most trustworthy and
    under what distribution shifts it fails.
\end{itemize}

\section{Project C --- Probing internal ``truthfulness'' representations}

\textbf{Goal.} Investigate whether a model's hidden states encode the truth of a
statement more reliably than its output, and whether such representations can be
read out without labels.

\textbf{Tasks.}
\begin{enumerate}[leftmargin=*]
  \item Train supervised linear probes on frozen activations to classify
    statements as true/false, sweeping over layers and token positions
    \citep{azaria2023internal}.
  \item Reproduce an unsupervised method that finds a truth direction via
    logical-consistency constraints (Contrast-Consistent Search) and compare it
    to the supervised probe \citep{burns2023discovering}.
  \item Measure the gap between what the probe ``believes'' and what the model
    \emph{says}: identify cases where the internal probe is correct while the
    generated answer is a hallucination, and quantify how often this occurs.
\end{enumerate}

\textbf{Deliverables.}
\begin{itemize}[leftmargin=*]
  \item Layerwise probe-accuracy curves plus a comparison of supervised vs.\
    unsupervised probes.
  \item A quantitative ``internal-knows-but-outputs-wrong'' statistic with
    examples.
  \item A critical appendix on probing pitfalls (spurious features, probe
    capacity, generalization across datasets).
\end{itemize}

\section{Project D --- Interventions: grounding and decoding}

\textbf{Goal.} Reduce hallucination via retrieval grounding and inference-time
decoding changes, and measure the trade-off against fluency and task accuracy.

\textbf{Tasks.}
\begin{enumerate}[leftmargin=*]
  \item Implement a retrieval-augmented pipeline and measure how grounding
    affects faithfulness and knowledge-gap errors \citep{lewis2020rag}.
  \item Implement a contrastive decoding intervention that amplifies factual
    signals by contrasting layers or model scales, e.g.\ DoLa
    \citep{chuang2024dola}, and compare against greedy and nucleus sampling.
  \item Combine the best intervention with the abstention policy from Project B
    and report the net effect on confident-error rate.
\end{enumerate}

\textbf{Deliverables.}
\begin{itemize}[leftmargin=*]
  \item An ablation table (grounding $\times$ decoding $\times$ abstention)
    reporting accuracy, faithfulness, and confident-error rate.
  \item A qualitative analysis of failure modes introduced by each intervention
    (e.g.\ over-copying from retrieved text, fluency degradation).
  \item A one-page recommendation on which combination best balances truthful
    abstention with usefulness.
\end{itemize}

\section*{Submission and assessment}

Each project should be submitted as a reproducible repository (fixed seeds,
environment file, and a \texttt{README} with exact commands), a 6--8 page report
with figures, and a 10-minute recorded demo. Grading weights clarity of the
experimental design, correctness of the calibration/statistical analysis, and
depth of the mechanistic interpretation. Negative results are acceptable when
accompanied by sound methodology.

\bibliographystyle{plainnat}
\begin{thebibliography}{10}

\bibitem{azaria2023internal}
A.~Azaria and T.~Mitchell.
\newblock The internal state of an LLM knows when it's lying.
\newblock In \emph{Findings of EMNLP}, 2023.

\bibitem{burns2023discovering}
C.~Burns, H.~Ye, D.~Klein, and J.~Steinhardt.
\newblock Discovering latent knowledge in language models without supervision.
\newblock In \emph{International Conference on Learning Representations (ICLR)}, 2023.

\bibitem{chuang2024dola}
Y.-S. Chuang, Y.~Xie, H.~Luo, Y.~Kim, J.~Glass, and P.~He.
\newblock DoLa: Decoding by contrasting layers improves factuality in large language models.
\newblock In \emph{International Conference on Learning Representations (ICLR)}, 2024.

\bibitem{kadavath2022know}
S.~Kadavath, T.~Conerly, A.~Askell, et~al.
\newblock Language models (mostly) know what they know.
\newblock \emph{arXiv preprint arXiv:2207.05221}, 2022.

\bibitem{kalai2025hallucinate}
A.~T. Kalai, O.~Nachum, S.~S. Vempala, and E.~Zhang.
\newblock Why language models hallucinate.
\newblock \emph{arXiv preprint arXiv:2509.04664}, 2025.

\bibitem{lewis2020rag}
P.~Lewis, E.~Perez, A.~Piktus, et~al.
\newblock Retrieval-augmented generation for knowledge-intensive NLP tasks.
\newblock In \emph{Advances in Neural Information Processing Systems (NeurIPS)}, 2020.

\bibitem{lin2022truthfulqa}
S.~Lin, J.~Hilton, and O.~Evans.
\newblock TruthfulQA: Measuring how models mimic human falsehoods.
\newblock In \emph{Proceedings of the Association for Computational Linguistics (ACL)}, 2022.

\bibitem{manakul2023selfcheckgpt}
P.~Manakul, A.~Liusie, and M.~J.~F. Gales.
\newblock SelfCheckGPT: Zero-resource black-box hallucination detection for generative large language models.
\newblock In \emph{Proceedings of EMNLP}, 2023.

\bibitem{maynez2020faithfulness}
J.~Maynez, S.~Narayan, B.~Bohnet, and R.~McDonald.
\newblock On faithfulness and factuality in abstractive summarization.
\newblock In \emph{Proceedings of the Association for Computational Linguistics (ACL)}, 2020.

\bibitem{openai2023gpt4}
OpenAI.
\newblock GPT-4 technical report.
\newblock \emph{arXiv preprint arXiv:2303.08774}, 2023.

\end{thebibliography}

\end{document}